% Notes of the Montreal seminar, 1970 (cote 129, typed part): archive cote 129, site page 49 (PDF page 50).
% typescript. Draft compared: none.
% AI-transcribed by Claude Opus 5.5 (claude-opus-5-5, adaptive thinking, effort high), Message Batches API, batch mtl129_full (msgbatch_01AGnDKHwy6ATCM7ehUa6263), 2026-11-13; not yet checked by eye.
\scanpage{129/49}
\mspage{\hand{II 21}}

5.3. \uline{Remarques}

5.3.1. Les groupes formels et les groupes de Barsotti-Tate sur $k$
peuvent être considérés comme limites inductives de p-groupes finis sur
$k$ . La théorie de Dieudonné permet donc de décrire ces groupes en termes
de limites projectives de $\mathcal{D}$-modules de longueur finie sur $W$ . Par exem-%
ple, on montrera (III, 5.6) que l'on a une équivalence de catégories en-%
tre la catégorie des groupes de Barsotti-Tate sur $k$ et la catégorie
des $\mathcal{D}$-modules qui sont libres et de type fini sur $W$ .

5.3.2. Pour un p-groupe $G$ fini sur $k$ , il existe un isomorphisme
canonique de $\mathcal{D}$-modules, fonctoriel en $G$ :
\[ D^*(G^*) \overset{\sim}{\to} \mathrm{Hom}_W(D^*(G),\ K/W) \ , \]
où $K/W$ désigne le module dualisant de $W$ (c.f. Gabriel et Demazure).
La structure de $\mathcal{D}$-module sur le membre de droite est définie comme suit :

Pour tout $\alpha \in \mathrm{Hom}_W(D^*(G),\ K/W)$ et tout $x \in D^*(G)$ , on pose
\[ F{\cdot}\alpha(x) = [\alpha(V(x))]^{\sigma} \]
\[ V{\cdot}\alpha(x) = [\alpha(F(x))]^{\sigma-1} \ . \]

6. \uline{Annexe = Construction du foncteur quasi inverse de} $D^*$

Il suffit de construire le foncteur quasi-inverse du foncteur
\diagram{
\[ D^* : \quad \left(\begin{array}{c} \text{groupes finis} \\ \text{unipotents} \end{array}\right)^{0} \longrightarrow \left(\begin{array}{l} \text{sous catégorie de } \underline{\mathrm{Mod}}_{\mathcal{D};\,W\ \hand{\ell.f.}} \text{ for-} \\ \text{mée des objets où } V \text{ est nilpotent} \end{array}\right) \]
}
défini en 4.3. On montrera que le foncteur que l'on va construire est un

\mspage{\hand{49}}
